\chapter{L18 -- Realizzazione dei sistemi, interconnessioni SISO, sistemi MIMO e forma minima}
\label{chap:L18}

\section{Richiamo iniziale: quando compare il termine \(t e^{-t}\)?}

Prima di entrare negli argomenti nuovi, conviene fissare bene un fatto che nei tuoi appunti compare come richiamo.

Se in una risposta temporale compare un termine del tipo
\[
t e^{-t},
\]
allora l'autovalore \(\lambda=-1\) \emph{non} è semplice nel senso dinamico: deve essere associato ad almeno un blocco di Jordan di ordine maggiore o uguale a \(2\).

Più precisamente:
\begin{itemize}
    \item se \(-1\) è autovalore semplice e \(A\) è diagonalizzabile in quel modo, allora i modi associati a \(-1\) sono solo multipli di \(e^{-t}\);
    \item il termine \(t e^{-t}\) compare solo se esiste una catena di Jordan di lunghezza almeno \(2\) associata a \(-1\);
    \item in generale, per un autovalore \(\lambda\), la dimensione del nucleo di \(A-\lambda I\) dà il numero di blocchi di Jordan associati a \(\lambda\), mentre la molteplicità algebrica dà la somma delle loro dimensioni.
\end{itemize}

Questa osservazione è importante perché, quando si studiano i modi di un sistema, non basta conoscere gli autovalori: bisogna anche capire la loro \emph{struttura di Jordan}.

\bigskip

\section{Forma minima di una realizzazione}

Sia
\[
G(s)=\frac{n(s)}{d(s)},
\qquad \deg d(s)=n.
\]

Una \emph{realizzazione minima} di \(G(s)\) è una realizzazione in forma di stato con:
\begin{itemize}
    \item numero minimo di stati;
    \item coppia \((A,B)\) completamente raggiungibile;
    \item coppia \((A,C)\) completamente osservabile.
\end{itemize}

Dal punto di vista strutturale, la forma minima si ottiene prendendo la decomposizione di Kalman e conservando solo il blocco contemporaneamente:
\begin{itemize}
    \item raggiungibile;
    \item osservabile.
\end{itemize}

Se indichiamo tale blocco con \((A_{RO},B_{RO},C_{RO})\), allora la funzione di trasferimento è

\[
G(s)=C_{RO}(sI-A_{RO})^{-1}B_{RO}+D.
\]

Questa è l'idea fondamentale:
\begin{quote}
\emph{la funzione di trasferimento ``vede'' solo la parte del sistema che è sia raggiungibile sia osservabile.}
\end{quote}

\bigskip

\section{Esercizio: equazione alle differenze e cancellazione polo-zero}

Consideriamo il sistema discreto descritto dall'equazione alle differenze
\[
y(t+2)+3y(t+1)+2y(t)=u(t+1)+2u(t),
\]
con condizioni iniziali nulle.

\subsection{Calcolo della funzione di trasferimento}

Applicando la trasformata \(Z\) con condizioni iniziali nulle si ottiene

\[
(z^2+3z+2)Y(z)=(z+2)U(z),
\]
quindi
\[
G(z)=\frac{Y(z)}{U(z)}=\frac{z+2}{z^2+3z+2}.
\]

Poiché
\[
z^2+3z+2=(z+1)(z+2),
\]
segue che
\[
G(z)=\frac{z+2}{(z+1)(z+2)}=\frac{1}{z+1}.
\]

Dunque c'è una cancellazione polo-zero nel punto \(z=-2\).

\subsection{Interpretazione della cancellazione}

La cancellazione dice che una realizzazione costruita \emph{prima} della riduzione può contenere uno stato in più rispetto al minimo necessario.

In particolare:
\begin{itemize}
    \item una realizzazione in forma canonica di controllo costruita dal denominatore non ridotto sarà completamente raggiungibile, ma non completamente osservabile;
    \item una realizzazione in forma canonica di osservazione costruita dal denominatore non ridotto sarà completamente osservabile, ma non completamente raggiungibile;
    \item se invece si realizza il sistema \emph{dopo} la cancellazione, si ottiene una realizzazione minima di ordine \(1\), che è sia raggiungibile sia osservabile.
\end{itemize}

\subsection{Una realizzazione minima}

Dalla forma ridotta
\[
G(z)=\frac{1}{z+1}
\]
si ottiene una realizzazione minima scalare, ad esempio

\[
x(t+1)=-x(t)+u(t),
\qquad
y(t)=x(t),
\]
cioè
\[
A=(-1),\qquad B=(1),\qquad C=(1),\qquad D=0.
\]

Essendo di ordine \(1\) e non avendo più cancellazioni, questa realizzazione è minima.

\bigskip

\section{Interconnessione di sistemi SISO}

Passiamo ora a un tema nuovo: come si realizza in forma di stato l'interconnessione di due sistemi.

\subsection{Connessione in serie}

Consideriamo due sistemi SISO \(\Sigma_1\) e \(\Sigma_2\), con funzioni di trasferimento
\[
G_1(s)=\frac{n_1(s)}{d_1(s)},
\qquad
G_2(s)=\frac{n_2(s)}{d_2(s)}.
\]

La connessione in serie è schematicamente

\[
u \longrightarrow \boxed{\Sigma_1} \longrightarrow \boxed{\Sigma_2} \longrightarrow y.
\]

Supponiamo che le realizzazioni siano

\[
\Sigma_1:
\begin{cases}
\dot x_1=A_1x_1+B_1u_1,\\
y_1=C_1x_1+D_1u_1,
\end{cases}
\qquad
\Sigma_2:
\begin{cases}
\dot x_2=A_2x_2+B_2u_2,\\
y_2=C_2x_2+D_2u_2.
\end{cases}
\]

Nella connessione in serie vale
\[
u_1=u,\qquad u_2=y_1,\qquad y=y_2.
\]

Quindi:
\[
u_2=C_1x_1+D_1u.
\]

Sostituendo nel secondo sistema:

\[
\dot x_2=A_2x_2+B_2(C_1x_1+D_1u)
      =A_2x_2+B_2C_1x_1+B_2D_1u.
\]

Inoltre
\[
y=y_2=C_2x_2+D_2u_2
     =C_2x_2+D_2C_1x_1+D_2D_1u.
\]

Se definiamo lo stato complessivo
\[
x=\begin{bmatrix}x_1\\x_2\end{bmatrix},
\]
si ottiene la realizzazione complessiva

\[
\begin{cases}
\dot x=
\begin{bmatrix}
A_1 & 0\\
B_2C_1 & A_2
\end{bmatrix}x+
\begin{bmatrix}
B_1\\
B_2D_1
\end{bmatrix}u,\\[1.2em]
y=
\begin{bmatrix}
D_2C_1 & C_2
\end{bmatrix}x+D_2D_1u.
\end{cases}
\]

Naturalmente, a livello di funzione di trasferimento, vale
\[
G(s)=G_2(s)G_1(s).
\]

\subsection{Cosa succede se ci sono cancellazioni?}

Se nella cascata si verifica una cancellazione polo-zero, la realizzazione ottenuta non è più minima.

Ci sono due casi concettualmente distinti.

\paragraph{Caso 1: uno zero di \(\Sigma_1\) coincide con un polo di \(\Sigma_2\).}

Allora il modo cancellato appartiene al secondo sistema, cioè a quello ``a valle''. In questo caso il sistema complessivo può perdere \emph{raggiungibilità}: esiste una dinamica interna di \(\Sigma_2\) che non riesce più a essere eccitata dall'ingresso esterno.

\paragraph{Caso 2: un polo di \(\Sigma_1\) coincide con uno zero di \(\Sigma_2\).}

In questo caso il modo cancellato appartiene al primo sistema, cioè a quello ``a monte''. Il sistema complessivo può perdere \emph{osservabilità}: esiste una dinamica interna che l'uscita finale non riesce più a vedere.

\medskip

In sintesi:
\begin{quote}
\emph{nelle interconnessioni in serie, le cancellazioni possono nascondere modi interni, rendendo il sistema complessivo non minimo.}
\end{quote}

\bigskip

\subsection{Connessione in parallelo}

La connessione in parallelo è schematicamente

\[
\begin{array}{c}
u \\
\downarrow
\end{array}
\qquad
\begin{array}{c}
\boxed{\Sigma_1}\\[0.4em]
\boxed{\Sigma_2}
\end{array}
\qquad
\longrightarrow\ y=y_1+y_2.
\]

Più precisamente:
\[
u_1=u,\qquad u_2=u,\qquad y=y_1+y_2.
\]

Se
\[
\Sigma_i:
\begin{cases}
\dot x_i=A_ix_i+B_iu,\\
y_i=C_ix_i+D_iu,
\end{cases}
\qquad i=1,2,
\]
e si pone
\[
x=\begin{bmatrix}x_1\\x_2\end{bmatrix},
\]
allora la realizzazione complessiva è

\[
\begin{cases}
\dot x=
\begin{bmatrix}
A_1 & 0\\
0 & A_2
\end{bmatrix}x+
\begin{bmatrix}
B_1\\
B_2
\end{bmatrix}u,\\[1.2em]
y=
\begin{bmatrix}
C_1 & C_2
\end{bmatrix}x+(D_1+D_2)u.
\end{cases}
\]

A livello di funzione di trasferimento:
\[
G(s)=G_1(s)+G_2(s)
=
\frac{n_1(s)}{d_1(s)}+\frac{n_2(s)}{d_2(s)}
=
\frac{n_1(s)d_2(s)+n_2(s)d_1(s)}{d_1(s)d_2(s)}.
\]

\subsection{Quando possono avvenire cancellazioni in parallelo?}

Nel collegamento in parallelo, eventuali cancellazioni nel sistema risultante possono comparire solo se i due sottosistemi hanno fattori dinamici comuni, cioè poli coincidenti nei denominatori.

In quel caso, il modo comune può sparire nella somma delle due risposte e il sistema complessivo non risulta più minimo.

Anche qui il messaggio è lo stesso:
\[
\text{cancellazione} \quad \Longrightarrow \quad \text{realizzazione non minima}.
\]

\bigskip

\section{Sistemi MIMO e matrice di trasferimento}

Passiamo ora ai sistemi con più ingressi e più uscite.

Se il sistema ha \(m\) ingressi e \(\ell\) uscite, la sua matrice di trasferimento è
\[
G(s)\in\mathbb R(s)^{\ell\times m},
\]
e si scrive come
\[
Y(s)=G(s)U(s),
\]
dove
\[
Y(s)=
\begin{bmatrix}
Y_1(s)\\ \vdots \\ Y_\ell(s)
\end{bmatrix},
\qquad
U(s)=
\begin{bmatrix}
U_1(s)\\ \vdots \\ U_m(s)
\end{bmatrix}.
\]

Gli elementi della matrice sono le funzioni di trasferimento ingresso-uscita:
\[
G(s)=
\begin{bmatrix}
g_{11}(s) & \cdots & g_{1m}(s)\\
\vdots & \ddots & \vdots\\
g_{\ell 1}(s) & \cdots & g_{\ell m}(s)
\end{bmatrix},
\]
e per ogni uscita vale
\[
Y_i(s)=g_{i1}(s)U_1(s)+g_{i2}(s)U_2(s)+\cdots+g_{im}(s)U_m(s).
\]

\subsection{Significato dei poli della matrice di trasferimento}

Per capire i poli di una matrice di trasferimento, una prima idea è spegnere tutti gli ingressi tranne uno.

Se fissiamo \(j\) e imponiamo
\[
U_k(s)=0 \qquad \forall k\neq j,
\]
allora
\[
Y(s)=G_{\bullet j}(s)\,U_j(s),
\]
dove \(G_{\bullet j}(s)\) è la \(j\)-esima colonna di \(G(s)\).

Questo significa che una colonna della matrice di trasferimento descrive \emph{tutti gli effetti del singolo ingresso \(u_j\) su tutte le uscite}. Analogamente, una riga descrive come una certa uscita dipende da tutti gli ingressi.

\subsection{Il termine diretto \(D\)}

Nel modello in forma di stato
\[
\begin{cases}
\dot x=Ax+Bu,\\
y=Cx+Du,
\end{cases}
\]
la matrice \(D\) raccoglie i termini di \emph{accoppiamento diretto} ingresso--uscita.

In una matrice di trasferimento MIMO, se un elemento \(g_{ij}(s)\) non è strettamente proprio, lo si può decomporre come
\[
g_{ij}(s)=d_{ij}+\text{parte strettamente propria},
\]
e il termine costante \(d_{ij}\) va inserito nella matrice \(D\).

\bigskip

\section{Realizzazione di una matrice di trasferimento lavorando per colonne}

Un metodo pratico per realizzare un sistema MIMO consiste nel lavorare \emph{per colonne}.

Supponiamo, ad esempio, di avere due ingressi e tre uscite, cioè una matrice di trasferimento \(G(s)\in\mathbb R(s)^{3\times 2}\). La scriviamo per colonne:
\[
G(s)=
\begin{bmatrix}
\vert & \vert\\
G^{(1)}(s) & G^{(2)}(s)\\
\vert & \vert
\end{bmatrix},
\]
dove ciascuna colonna \(G^{(j)}(s)\) contiene l'effetto dell'ingresso \(u_j\) su tutte le uscite.

Per ogni colonna:
\begin{enumerate}
    \item si cerca un denominatore comune;
    \item si costruisce una realizzazione SISO/MISO in forma canonica di controllo;
    \item si accoppiano poi i vari blocchi in modo disaccoppiato sugli ingressi.
\end{enumerate}

\subsection{Struttura generale ottenuta per colonne}

Se dalla colonna \(j\)-esima si ricava una realizzazione
\[
(A^{(j)},B^{(j)},C^{(j)}),
\]
allora una realizzazione complessiva si può scrivere come

\[
A=
\begin{bmatrix}
A^{(1)} & 0 & \cdots & 0\\
0 & A^{(2)} & \cdots & 0\\
\vdots & \vdots & \ddots & \vdots\\
0 & 0 & \cdots & A^{(m)}
\end{bmatrix},
\qquad
B=
\begin{bmatrix}
B^{(1)} & 0 & \cdots & 0\\
0 & B^{(2)} & \cdots & 0\\
\vdots & \vdots & \ddots & \vdots\\
0 & 0 & \cdots & B^{(m)}
\end{bmatrix},
\]
e
\[
C=
\begin{bmatrix}
C^{(1)} & C^{(2)} & \cdots & C^{(m)}
\end{bmatrix}.
\]

Questa costruzione è molto utile perché permette di realizzare la matrice di trasferimento in modo operativo. Tuttavia:
\begin{itemize}
    \item la realizzazione ottenuta \emph{non è in generale minima};
    \item potrebbero esserci cancellazioni tra colonne;
    \item spesso è necessario ridurre poi la realizzazione con strumenti strutturali.
\end{itemize}

\bigskip

\section{Metodo duale: lavorare per righe}

Esiste naturalmente il metodo duale, che consiste nel lavorare per \emph{righe}.

In questo caso si guarda ogni uscita separatamente:
\[
Y_i(s)=g_{i1}(s)U_1(s)+\cdots+g_{im}(s)U_m(s).
\]

Dal punto di vista realizzativo:
\begin{itemize}
    \item lavorare per colonne enfatizza il ruolo degli ingressi;
    \item lavorare per righe enfatizza il ruolo delle uscite;
    \item i due metodi sono duali.
\end{itemize}

Non c'è una regola assoluta per scegliere tra i due:
\begin{quote}
\emph{si sceglie di solito il metodo più conveniente, cioè quello che porta a sottosistemi di ordine più basso o con strutture più semplici.}
\end{quote}

\bigskip

\section{Come si individuano poli e zeri di una matrice di trasferimento?}

Per sistemi MIMO la nozione di poli e zeri non si legge più guardando semplicemente ogni singolo elemento della matrice. Serve uno strumento più robusto: la \emph{forma di Smith} di una matrice polinomiale.

\section{Forma di Smith di una matrice polinomiale}

Sia
\[
P(s)\in \mathbb R[s]^{p\times m}
\]
una matrice polinomiale.

Allora esistono matrici unimodulari \(U_1(s)\) e \(U_2(s)\) tali che
\[
P(s)=U_1(s)\,\Lambda(s)\,U_2(s),
\]
dove
\[
\Lambda(s)=
\operatorname{diag}\bigl(d_1(s),d_2(s),\dots,d_r(s),0,\dots,0\bigr),
\]
con:
\begin{itemize}
    \item \(r=\operatorname{rk}(P(s))\);
    \item i polinomi \(d_i(s)\) sono monici;
    \item \(d_i(s)\) divide \(d_{i+1}(s)\) per ogni \(i\).
\end{itemize}

Questa è la \emph{forma di Smith}.

\subsection{Come si costruiscono i polinomi invarianti}

Si definiscono:
\[
D_0(s)=1,
\]
e, per \(i=1,\dots,r\),
\[
D_i(s)=\gcd\{\text{tutti i minori di ordine } i \text{ di } P(s)\},
\]
presi monici.

Poi si pongono
\[
d_i(s)=\frac{D_i(s)}{D_{i-1}(s)}.
\]

I polinomi \(d_i(s)\) sono proprio gli \emph{invariant factors} della matrice.

\bigskip

\section{Esempio di forma di Smith}

Consideriamo
\[
P(s)=
\begin{bmatrix}
1 & -1\\
s^2+s-4 & 2s^2-s-8\\
s^2-4 & 2s^2-8
\end{bmatrix}.
\]

\subsection{Calcolo di \(D_1(s)\)}

Il massimo comun divisore di tutti gli elementi è
\[
D_1(s)=1.
\]

\subsection{Calcolo di \(D_2(s)\)}

Occorre calcolare i minori \(2\times 2\).

\paragraph{Minore delle righe 1 e 2:}
\[
\det
\begin{bmatrix}
1 & -1\\
s^2+s-4 & 2s^2-s-8
\end{bmatrix}
=
2s^2-s-8+s^2+s-4
=
3s^2-12
=
3(s^2-4).
\]

\paragraph{Minore delle righe 1 e 3:}
\[
\det
\begin{bmatrix}
1 & -1\\
s^2-4 & 2s^2-8
\end{bmatrix}
=
2s^2-8+s^2-4
=
3s^2-12
=
3(s^2-4).
\]

\paragraph{Minore delle righe 2 e 3:}
\[
\det
\begin{bmatrix}
s^2+s-4 & 2s^2-s-8\\
s^2-4 & 2s^2-8
\end{bmatrix}.
\]

Sviluppando:
\[
(s^2+s-4)(2s^2-8)-(2s^2-s-8)(s^2-4)
=
3s(s^2-4).
\]

Il massimo comun divisore monico di questi minori è quindi
\[
D_2(s)=s^2-4.
\]

\subsection{Polinomi invarianti}

Poiché
\[
D_0(s)=1,\qquad D_1(s)=1,\qquad D_2(s)=s^2-4,
\]
si ottiene
\[
d_1(s)=\frac{D_1(s)}{D_0(s)}=1,
\qquad
d_2(s)=\frac{D_2(s)}{D_1(s)}=s^2-4.
\]

Quindi la forma di Smith è

\[
\Lambda(s)=
\begin{bmatrix}
1 & 0\\
0 & s^2-4\\
0 & 0
\end{bmatrix}.
\]

\medskip

Questo esempio è importante perché mostra che, in un sistema MIMO, i poli e gli zeri strutturali non si leggono banalmente elemento per elemento, ma emergono dai fattori invarianti.

\bigskip

\section{Messaggio concettuale della lezione}

Possiamo riassumere l'intera lezione così.

\begin{enumerate}
    \item Una funzione di trasferimento può essere realizzata in molti modi diversi, ma la \emph{forma minima} è quella con il minor numero di stati ed è sia completamente raggiungibile sia completamente osservabile.
    \item Nelle interconnessioni in serie e in parallelo possono comparire cancellazioni interne, che rendono la realizzazione complessiva non minima.
    \item Nei sistemi MIMO la matrice di trasferimento va studiata come oggetto unico, non guardando soltanto i singoli elementi.
    \item Per realizzare un sistema MIMO si può lavorare per colonne o per righe, scegliendo la via più conveniente.
    \item Per studiare poli e zeri strutturali di una matrice di trasferimento si usa la forma di Smith della matrice polinomiale associata.
\end{enumerate}

\bigskip

\section*{Promemoria operativo}

Quando ti danno una matrice di trasferimento MIMO e ti chiedono una realizzazione, la strategia pratica è:

\begin{enumerate}
    \item separare l'eventuale parte diretta \(D\), cioè i termini non strettamente propri;
    \item decidere se conviene lavorare per colonne o per righe;
    \item per ciascun blocco trovato, costruire una realizzazione canonica del denominatore comune;
    \item assemblare i blocchi in una realizzazione complessiva;
    \item verificare se la realizzazione è minima oppure se restano modi cancellati;
    \item se serve, usare la decomposizione di Kalman o la forma di Smith per ridurre il sistema.
\end{enumerate}